\subsection{Algebraic extensions of a finite field}

PROPOSITION 3. --- Let K be a finite field of q elements, Ω an algebraically closed extension of K and m an integer ≥1.

\begin{enumerate}
\def\labelenumi{\alph{enumi})}
\item
  There exists a unique subextension \(\mathbf{K}_m\) of \(\Omega\) which is of degree \(m\) over \(\mathbf{K}\) .
\item
  The field \(K_{m}\) has \(q^{m}\) elements and is the set of fixed points of the automorphism \(x \mapsto x^{q^{m}}\) of \(\Omega\) .
\item
  We have \(\mathbf{K}_m = \mathbf{K}(\zeta)\) for every generator \(\zeta\) of the cyclic group \(\mathbf{K}_m^*\) .
\end{enumerate}

Let \(p\) be the characteristic of \(\mathbf{K}\) and \(f\) the degree of \(\mathbf{K}\) over \(\mathbf{F}_p\) . We have \(q^m = p^{fm}\) and the mapping \(x \mapsto x^{q^m}\) is thus an automorphism of the perfect field \(\Omega\) (V, p.~7, Prop. 4). Therefore the set \(\mathbf{K}_m\) of roots of the polynomial \(X^{q^m} - X\) of \(\mathbf{K}[\mathbf{X}]\) is a subfield of \(\Omega\) . Since the derivative of \(X^{q^m} - X\) is equal to -1, all the roots of this polynomial are simple (IV, p.~17, Prop. 7) and so \(\mathbf{K}_m\) has \(q^m\) elements. It follows that \([\mathbf{K}_m : \mathbf{K}] = m\) .

Now let L be a subextension of \(\Omega\) , of degree \(m\) over K. As vector space over K, L is isomorphic to \(\mathbf{K}^m\) and so has \(q^m\) elements. We thus have \(x^{q^m} = x\) for all \(x \in L\) (Prop. 2), whence \(\mathrm{L} \subset \mathrm{K}_m\) . Since \([\mathrm{L} : \mathrm{K}] = [\mathrm{K}_m : \mathrm{K}] = m\) , we finally have \(\mathrm{L} = \mathrm{K}_m\) .

We have thus proved Assertions \(a)\) and \(b)\) , and \(c)\) is trivial.

COROLLARY. --- Let K be a finite field and \(\Omega\) an algebraically closed extension of K. The relative algebraic closure \(\bar{K}\) of K in \(\Omega\) consists of 0 and roots of unity and is an algebraic closure of K.

We know \(\bar{K}\) to be an algebraic closure of K (V, p.~22, Example 2) and it is clear that every root of unity in \(\Omega\) belongs to \(\bar{K}\) . Further, given \(x \neq 0\) in \(\bar{K}\) , of degree m over K, if K has q elements, then the field \(\mathbf{K}(x)\) has \(q^{m}\) , whence \(x^{q^{m}-1} = 1\) , and so x is a root of unity in \(\Omega\) .

Let \(p\) be a prime number and let \(\mathbf{F}_p = \mathbf{Z} / p\mathbf{Z}\) be the field of \(p\) elements. We choose an algebraic closure \(\Omega\) of \(\mathbf{F}_p\) , whose existence follows from Steinitz's theorem (V, p.~23, Th. 2). Let \(f\) be a positive integer and \(q = p^f\) . By Prop. 3 there exists a unique subfield of \(\Omega\) which is of degree \(f\) over \(\mathbf{F}_p\) ; we shall denote it by \(\mathbf{F}_q(\Omega)\) or by abuse of notation \(\mathbf{F}_q\) . It is the unique subextension of \(\Omega\) which has degree \(f\) over \(\mathbf{F}_p\) . It is the unique subfield of \(\Omega\) of cardinal \(q\) , and every field of cardinal \(q\) is isomorphic (not canonically) to \(\mathbf{F}_q\) (Cor. of Prop. 2). We note that \(\mathbf{F}_q\) consists of those \(x\) in \(\Omega\) for which \(x^q = x\) and that \(\mathbf{F}_q \subset \mathbf{F}_{q'}\) if and only if \(q'\) is a power of \(q\) .

PROPOSITION 4. --- Let K be a finite field of q elements and \(K_{m}\) an extension of finite degree m of K.

\begin{enumerate}
\def\labelenumi{\alph{enumi})}
\item
  The field \(K_{m}\) is a Galois extension of K whose Galois group is the cyclic group of order m generated by the automorphism \(\sigma_{q}: x \mapsto x^{q}\) .
\item
  For each \(x \in \mathbf{K}_m\) the norm of \(x\) with respect to \(\mathbf{K}\) is equal to \(x^{(q^m - 1) / (q - 1)}\) .
\item
  Every element of \(\mathbf{K}\) is the trace (resp. norm) of an element of \(\mathbf{K}_m\) .
\end{enumerate}

Let \(\Gamma\) be the cyclic group of automorphisms of \(\mathbf{K}_m\) generated by \(\sigma_q\) . The field of invariants of \(\Gamma\) consists of the elements \(x\) of \(\mathbf{K}_m\) such that \(x^q = x\) , hence is equal to \(\mathbf{K}\) . Therefore \(\mathbf{K}_m\) is a Galois extension of \(\mathbf{K}\) with Galois group \(\Gamma\) , and this latter has order equal to \([\mathbf{K}_m: \mathbf{K}] = m\) (V, p.~66, Th. 3). Thus \(a)\) follows.

We have \(\Gamma = \{1, \sigma_q, \sigma_q^2, \ldots, \sigma_q^{m-1}\}\) ; the norm of an element \(x\) of \(\mathbf{K}_m\) with respect to \(\mathbf{K}\) is thus \(\mathbf{N}(x) = \prod_{i=0}^{m-1} \sigma_q^i(x) = x^{1+q+\cdots+q^{m-1}}\) and we have \(1 + q + \cdots + q^{m-1} = \frac{q^m - 1}{q - 1}\) . This proves \(b\) ). Let \(\zeta\) be a generator of the cyclic group \(\mathbf{K}_m^*\) ; the image of the norm \(\mathbf{N}: \mathbf{K}_M^* \to \mathbf{K}^*\) is the cyclic subgroup of \(\mathbf{K}_i^*\) generated by the element \(\xi = \mathbf{N}(\zeta) = \zeta^{(q^m - 1)/(q - 1)}\) ; since \(\zeta\) is of order \(q^m - 1\) , \(\xi\) is of order \(q - 1\) , and so generates \(\mathbf{K}^*\) . This proves that every non-zero element of K is norm of a non-zero element of \(K_{m}\) ; moreover, we have \(0 = \mathbf{N}(0)\) .

Finally since \(K_{m}\) is a separable algebraic extension of K, the trace is a non-zero linear form on the vector space \(K_{m}\) over K (V, p.~49, Cor.) ; so every element of K is the trace of an element of \(K_{m}\) .
% semantic-fragments: legacy-input-seam
\relax{}
